\section*{Bab \ref{ch:b1:elimination} --- Eliminasi \texorpdfstring{$\beta$}{Beta}}

\begin{solution}{exo:b1:elimination:1}
2-Bromobutana: but-1-ena, \cip{E}- dan \cip{Z}-but-2-ena; yang utama adalah
\cip{E}-but-2-ena. 2-Bromo-2-metilbutana: 2-metilbut-2-ena (utama,
trisubstitusi) dan 2-metilbut-1-ena. Bromosikloheksana: hanya
sikloheksena.
\end{solution}

\begin{solution}{exo:b1:elimination:2}
E1: $v = k[\mathrm{RBr}]$; E2: $v = k[\mathrm{RBr}][\ce{EtO-}]$. Ukurlah
\omterm{def:b1:rate-laws:initial-rate}{laju awal} pembentukan alkena pada beberapa konsentrasi etoksida: laju itu
bertambah sebanding untuk E2, dan tidak berubah untuk E1.
\end{solution}

\begin{solution}{exo:b1:elimination:3}
Sepanjang C2--C3, Br pada C2 dapat anti terhadap salah satu hidrogen C3
atau terhadap gugus metil C4: kedua \omterm{def:b1:stereochemistry-in-depth:isomers}{konformasi} dengan H anti terhadap Br
dapat bereaksi melalui E2 (menghasilkan alkena \cip{E} dan \cip{Z});
\omterm{def:b1:stereochemistry-in-depth:isomers}{konformasi} ketiga tidak dapat mengalami eliminasi ke arah C3.
\end{solution}

\begin{solution}{exo:b1:elimination:4}
SN2 (primer, hidroksida, suhu rendah); E2 (tersier, \omterm{def:b1:predominant-reaction:strong-acid}{basa kuat} yang
meruah); SN2 (iodida, \omterm{def:b1:electronic-effects:nucleophile}{nukleofil} yang baik dan \omterm{def:b1:predominant-reaction:strong-acid}{basa lemah}, \omterm{def:b1:intermolecular-forces-solvents:solvent-classes}{pelarut
aprotik}); SN1 dan E1, dengan E1 bertambah jika dipanaskan.
\end{solution}

\begin{solution}{exo:b1:elimination:5}
C3 membawa satu hidrogen. Pada setiap \omterm{def:b1:stereochemistry-in-depth:enantiomers}{diastereomer}, menempatkan H itu anti
terhadap Br pada C2 menetapkan posisi gugus-gugus lainnya: metil C2 dan
etil C3 berada di sisi yang sama pada \omterm{def:b1:stereochemistry-in-depth:enantiomers}{diastereomer} yang satu dan di sisi
yang berlawanan pada \omterm{def:b1:stereochemistry-in-depth:enantiomers}{diastereomer} yang lain. Yang satu menghasilkan
\cip{E}-, yang lain \cip{Z}-3-metilpent-2-ena: sebuah \omterm{def:b1:nucleophilic-substitution:stereospecific}{reaksi
stereospesifik}.
\end{solution}

\begin{solution}{exo:b1:elimination:6}
OH terprotonasi, \ce{R-OH2+}; air lepas, menghasilkan \omterm{def:b1:electronic-effects:intermediates}{karbokation} tersier
\ce{(CH3)2C+-CH2CH3}; sebuah basa (air, \ce{HSO4-}) mengambil proton dari
karbon tetangganya. Produk utama (Zaitsev): 2-metilbut-2-ena.
\end{solution}

\begin{solution}{exo:b1:elimination:7}
Gugus tert-butil mengunci \omterm{def:b1:stereochemistry-in-depth:conformations}{kursi} dengan gugus itu \omterm{def:b1:stereochemistry-in-depth:conformations}{ekuatorial}. Pada isomer
\textit{trans}, brom lalu juga \omterm{def:b1:stereochemistry-in-depth:conformations}{ekuatorial}, tidak pernah anti terhadap
C--H: E2 harus menunggu \omterm{def:b1:stereochemistry-in-depth:conformations}{kursi} terbalik yang sangat tidak menguntungkan.
Pada isomer \textit{cis} brom \omterm{def:b1:stereochemistry-in-depth:conformations}{aksial}, dengan dua hidrogen trans-diaksial:
E2 berlangsung cepat.
\end{solution}

\begin{solution}{exo:b1:elimination:8}
Etoksida, yang kecil, mengambil hidrogen yang menghasilkan alkena yang
lebih tersubstitusi (Zaitsev). tert-Butoksida, yang meruah, lebih mudah
mencapai hidrogen-hidrogen gugus metil yang terbuka daripada hidrogen
\ce{CH2} di bagian dalam: alkena yang kurang tersubstitusi.
\end{solution}

\begin{solution}{exo:b1:elimination:9}
Bromometana tidak mempunyai karbon $\beta$. Pada 1-bromo-2,2-dimetilpropana
karbon $\beta$-nya kuaterner dan tidak membawa hidrogen.
\end{solution}

\begin{solution}{exo:b1:elimination:10}
Kedua produk memerlukan \omterm{def:b1:electronic-effects:intermediates}{karbokation}, yang terbentuk pada tahap yang lambat
dengan laju $k_1[\mathrm{RBr}]$. \omterm{def:b1:electronic-effects:intermediates}{Karbokation} itu lalu terbagi antara
penangkapan oleh etanol (\omterm{def:b1:rate-laws:order}{tetapan laju} $k_S$) dan pelepasan proton
(\omterm{def:b1:rate-laws:order}{tetapan laju} $k_E$), keduanya berorde satu semu terhadap pelarut: fraksi
alkenanya $k_E/(k_E + k_S)$, tidak bergantung pada konsentrasi \omterm{def:b1:nucleophilic-substitution:substitution}{substrat}.
\end{solution}

\begin{solution}{exo:b1:elimination:11}
Mulai dari \omterm{def:b1:stereochemistry-in-depth:isomers}{konformasi} \omterm{def:b1:stereochemistry-in-depth:enantiomers}{senyawa meso} dengan kedua Br anti (kedua fenil lalu
anti, kedua H anti), putarlah C1 sebesar $120^\circ$ untuk menempatkan
H-nya anti terhadap Br pada C2. Kedua gugus fenil lalu berada di sisi
yang sama terhadap sumbu H--Br: keduanya berakhir \textit{cis}. Pada C1,
Br berperingkat lebih tinggi daripada Ph; pada C2, Ph lebih tinggi
daripada H; Br dan fenil C2 bersifat \textit{trans}:
\cip{E}-1-bromo-1,2-difeniletena.
\end{solution}

\begin{solution}{exo:b1:elimination:12}
Hanya \omterm{def:b1:stereochemistry-in-depth:conformations}{kursi} dengan \omterm{def:b1:electronic-effects:nucleophile}{gugus pergi} \omterm{def:b1:stereochemistry-in-depth:conformations}{aksial} yang bereaksi: $v = k\,x\,[\mathrm{RY}][\mathrm{B}]$,
dengan $x$ fraksi \omterm{def:b1:stereochemistry-in-depth:conformations}{kursi} itu. Setiap gugus alkil \omterm{def:b1:stereochemistry-in-depth:conformations}{aksial} yang dituntutnya
memerlukan beberapa kJ/mol (\qty{5.7}{kJ/mol} untuk metil), yang membagi
$x$ sekitar sepuluh kali per gugus pada suhu kamar: dengan dua atau tiga
gugus semacam itu, $x$ turun menjadi $10^{-2}$ atau $10^{-3}$ dan
reaksinya sebanyak itu pula lebih lambat.
\end{solution}

\begin{solution}{pb:b1:elimination:1}
\textbf{1.} Sebuah \omterm{def:b1:stereochemistry-in-depth:conformations}{kursi} dengan Cl, isopropil, dan metil semuanya
\omterm{def:b1:stereochemistry-in-depth:conformations}{ekuatorial}.
\textbf{2.} Isopropil dan metil \omterm{def:b1:stereochemistry-in-depth:conformations}{ekuatorial}, Cl \omterm{def:b1:stereochemistry-in-depth:conformations}{aksial}.
\textbf{3.} \omterm{def:b1:stereochemistry-in-depth:enantiomers}{Diastereomer}: keduanya hanya berbeda pada C1.
\textbf{4.} \omterm{def:b1:electronic-effects:nucleophile}{Gugus pergi} dan hidrogen $\beta$ trans-diaksial.
\textbf{5.} Hanya pada \omterm{def:b1:stereochemistry-in-depth:conformations}{kursi} terbalik, dengan Cl, isopropil, dan metil
semuanya \omterm{def:b1:stereochemistry-in-depth:conformations}{aksial}.
\textbf{6.} Ketiga gugus \omterm{def:b1:stereochemistry-in-depth:conformations}{aksial}: metil saja memerlukan sekitar
\qty{5.7}{kJ/mol}, gugus isopropil yang lebih besar lebih banyak, klor
sedikit: seluruhnya sekitar \qty{13}{kJ/mol}, sehingga fraksinya sekitar
$\exp(-13000/2479) = \num{5e-3}$, nilai yang diambil dalam soal.
\textbf{7.} C2 (satu H, karena membawa gugus isopropil) dan C6 (dua H).
\textbf{8.} Dengan Cl \omterm{def:b1:stereochemistry-in-depth:conformations}{aksial} pada C1, H \omterm{def:b1:stereochemistry-in-depth:conformations}{aksial} C2 (karena isopropilnya
\omterm{def:b1:stereochemistry-in-depth:conformations}{ekuatorial}) dan H \omterm{def:b1:stereochemistry-in-depth:conformations}{aksial} C6: dua.
\textbf{9.} Pada \omterm{def:b1:stereochemistry-in-depth:conformations}{kursi} triaksial gugus isopropil menempati posisi aksial
C2, sehingga H-nya \omterm{def:b1:stereochemistry-in-depth:conformations}{ekuatorial}: tidak anti. C6 mempunyai H \omterm{def:b1:stereochemistry-in-depth:conformations}{aksial}: anti.
\textbf{10.} Neomentil klorida dua, mentil klorida satu.
\textbf{11.} Sepanjang C1--C6: Cl (\omterm{def:b1:stereochemistry-in-depth:conformations}{aksial}, ke atas) pada C1 berseberangan
dengan H \omterm{def:b1:stereochemistry-in-depth:conformations}{aksial} (ke bawah) pada C6; ikatan-ikatan cincin dan H
\omterm{def:b1:stereochemistry-in-depth:conformations}{ekuatorial} goyang di antaranya.
\textbf{12.} Pada \omterm{def:b1:stereochemistry-in-depth:conformations}{kursi} tidak ada ikatan C--H yang \omterm{def:b1:elimination:periplanar}{sin-periplanar}
terhadap ikatan C--Cl; memaksakannya berarti cincin yang eklips dan
tegang.
\textbf{13.} Melepaskan H C2: 4-metil-1-(propan-2-il)sikloheks-1-ena
(trisubstitusi); melepaskan H C6: 3-metil-6-(propan-2-il)sikloheks-1-ena
(disubstitusi).
\textbf{14.} Alkena trisubstitusi, yang lebih tersubstitusi dan lebih
stabil.
\textbf{15.} Kedua alkena, \qty{78}{\percent} trisubstitusi: sesuai.
\textbf{16.} Hanya alkena disubstitusi: kebalikan ramalan Zaitsev, yang
dipaksakan oleh satu-satunya hidrogen anti yang tersedia.
\textbf{17.} Keduanya stereospesifik (syarat anti); neomentil klorida
bereaksi secara regioselektif (78/22); mentil klorida menghasilkan satu
\omterm{def:b1:stereochemistry-in-depth:isomers}{isomer struktur} saja.
\textbf{18.} Karbon yang membawa Cl bersifat sekunder dan diapit gugus
isopropil: padat bagi SN2, sedangkan etoksida adalah \omterm{def:b1:predominant-reaction:strong-acid}{basa kuat}.
\textbf{19.} $v \propto x\,n_{\mathrm{H}}$: fraksi \omterm{def:b1:stereochemistry-in-depth:conformations}{kursi} reaktif dikalikan
dengan banyaknya hidrogen anti.
\textbf{20.} $0.95 \times 2 = 1.90$.
\textbf{21.} $\num{5.0e-3} \times 1 = \num{5.0e-3}$.
\textbf{22.} \textbf{$k(\text{neomentil})/k(\text{mentil}) = 1.90/\num{5.0e-3}
= 380$}.
\end{solution}
